% titel: Komet Hale-Bopp: Keplersche Gesetze
% skript: Kap. Gravitation, Keplersche Gesetze

Beim Kometen Hale-Bopp (sehr langgestreckte Ellipsenbahn um die Sonne) sind
die drei Keplerschen Gesetze unmittelbar anwendbar: das dritte liefert aus der
Umlaufzeit die grosse Halbachse, das zweite (Drehimpulserhaltung) das
Geschwindigkeitsverhältnis in Perihel und Aphel.

\begin{werkzeugkasten}[Werkzeuge]
\textbf{1.\ Kepler:} Planeten/Kometen laufen auf Ellipsen, die Sonne in einem
Brennpunkt.

\textbf{2.\ Kepler / Drehimpulserhaltung:} In gleichen Zeiten überstreicht der
Fahrstrahl gleiche Flächen. Im Perihel (Abstand $r_P$) und Aphel ($r_A$)
-- wo die Geschwindigkeit senkrecht zum Fahrstrahl steht -- folgt
\begin{equation}
  L = m\, r_P\, v_P = m\, r_A\, v_A
  \;\Rightarrow\;
  \frac{v_P}{v_A} = \frac{r_A}{r_P}\;.
\end{equation}

\textbf{3.\ Kepler:} $T^2 \propto a^3$ ($a$ grosse Halbachse):
\begin{equation}
  \frac{T^2}{a^3} = \frac{4\pi^2}{G M_\odot}
  \;(\text{einerlei Zentralmasse}).
\end{equation}
Für Bahnen um die Sonne gilt in \emph{Jahren} und \emph{AE} exakt $T^2=a^3$.

\textbf{Potentielle Energie} im Gravitationsfeld der Sonne:
\begin{equation}
  E_\mathrm{pot} = -G\frac{M_\odot\, m}{r}\;.
\end{equation}

\textbf{AE} (Astronomische Einheit) $=\qty{1.5e8}{km}$; $G=\qty{6.674e-11}{m^3/(kg\,s^2)}$,
$M_\odot=\qty{1.989e30}{kg}$.
\end{werkzeugkasten}

\begin{werkzeugkasten}[Vorgehen in allgemeiner Form]
Aus $T$ die Halbachse $a$ (Teil a), daraus über $r_A = 2a-r_P$ das Aphel (b).
Geschwindigkeitsverhältnis aus dem Drehimpuls (c). Die Zeit Perihel$\to$Aphel
näherungsweise als halber Ellipsenumlauf bzw. Strecke durch mittlere
Geschwindigkeit (d). Potentielle Energien an beiden Bahnpunkten (e).
\end{werkzeugkasten}

\begin{achtung}[Typische Fallen]
\begin{itemize}
  \item $r_P$ und $r_A$ sind Abstände von der \emph{Sonne}, und
       $2a=r_P+r_A$.
  \item Das Produkt aus Geschwindigkeit und Abstand ist konstant (nicht Geschwindigkeit
       allein); das Verhältnis dreht sich um.
  \item AE-Einheiten: bei $T^2=a^3$ muss $T$ in Jahren, $a$ in AE stehen.
\end{itemize}
\end{achtung}
